
% !TeX program = lualatex
% !TeX encoding = utf8
% !TeX spellcheck = uk_UA
% !BIB program = bibtex8


\documentclass[]{article}
\usepackage[%
	a4paper,%
	footskip=1cm,%
	headsep=0.3cm,%
	top=1cm, %поле сверху
	bottom=2cm, %поле снизу
	left=2cm, %поле ліворуч
	right=2cm, %поле праворуч
    ]{geometry}

\usepackage[dvipsnames]{xcolor}
\usepackage[fontsize=12pt]{fontsize}
\usepackage{fontspec}
\setsansfont{CMU Sans Serif}%{Arial}
\setmainfont{CMU Serif}%{Times New Roman}
\setmonofont{CMU Typewriter Text}%{Consolas}
\defaultfontfeatures{Ligatures={TeX}}
\usepackage[math-style=ISO]{unicode-math}
\setmathfont{STIX Two Math}
\usepackage{amsmath}
\usepackage[english, russian, ukrainian]{babel}

\usepackage{enumitem}
\usepackage{tikz}

\allowdisplaybreaks
\usepackage{cancel}

\def\ket#1{\left| #1 \right\rangle }
\def\bra#1{\left\langle #1 \right| }
\def\bracket#1#2{\left\langle #1\left| #2 \right. \right\rangle }
\def\opbracket#1#2#3{\left\langle #1\left| \textcolor{OliveGreen}{#2}\right|  #3 \right\rangle}
\def\eff{e\kern-2.4pt{f\kern-2.65pt f}}
\def\vxi{\vec{\xi}}
\def\vr{\vec{r}}
\let\vphi\varphi
\let\ve\varepsilon

\def\f{(\textcolor{red}{1})}
\def\s{(\textcolor{blue}{2})}

\def\F{\textcolor{red}{\mathbf{1}}}
\def\S{\textcolor{blue}{\mathbf{2}}}

\def\oa{\textcolor{red}{ \vphi_{\mathbf{1}}} }
\def\ob{\textcolor{blue}{ \vphi_{\mathbf{2}}} }
\def\cPhi{\oa\f \ob\s - \oa\s \ob\f} % лінійна Комбінація
\def\h{\hat{h}}
\def\ha{\h\f}
\def\hb{\h\s}
\def\V{\frac{1}{r_{12}}}
\def\a{\oa\f \ob\s}
\def\b{\ob\f \oa\s}
\def\dt{\,.\,}
%\def\sp{\textvisiblespace}

\def\hr{\hrulefill\\*[10pt]}
\def\hhr{\hrulefill\\[-0.9\baselineskip]\hrulefill\\*[10pt]}

\title{Метод Хартрі--Фока для двоелектронної системи (атом гелію)}
\author{}
\date{}

\begin{document}
\maketitle

\section{Постановка задачі}

Розглянемо систему двох електронів, що рухаються в полі нерухомого ядра (наприклад, атом гелію або гелієподібний іон). Наближення Хартрі--Фока полягає в тому, що точна багатоелектронна хвильова функція замінюється \emph{однією} антисиметризованою комбінацією одноелектронних функцій (спін-орбіталей) --- визначником Слейтера. Далі, варіюючи енергію за цими одноелектронними функціями при додатковій умові їх ортонормованості, отримують рівняння, яким ці функції повинні задовольняти. Мета цього конспекту --- послідовно, крок за кроком, вивести канонічні рівняння Хартрі--Фока для системи двох електронів та явний вигляд кулонівського й обмінного інтегралів.

Через усе виведення позначення $\f$ та $\s$ (виділені кольором) вказують, від координат \emph{якого} електрона --- першого чи другого --- залежить вираз; це суто позначення, а не множники.

\section{Хвильова функція системи}

Електрони --- тотожні ферміони, тому повна хвильова функція системи має бути антисиметричною відносно перестановки будь-якої пари електронів (принцип Паулі). Якщо електрон 1 займає одноелектронний стан $\oa$, а електрон 2 --- стан $\ob$, то єдина антисиметрична комбінація, яку можна з них побудувати, --- це визначник Слейтера $2\times2$:

\begin{equation*}
	\Phi(1, 2) ={1}/{\sqrt2}\left[\cPhi\right]
\end{equation*}

Множник $1/\sqrt2$ забезпечує нормування $\langle\Phi|\Phi\rangle=1$ за умови, що самі орбіталі $\oa$, $\ob$ ортонормовані. Легко перевірити, що перестановка $1\leftrightarrow2$ змінює знак усього виразу: $\Phi(2,1)=-\Phi(1,2)$, як і належить фермі-статистиці. Якщо покласти $\oa=\ob$, визначник обертається в нуль --- це і є принцип заборони Паулі для двох електронів у стані з однаковими квантовими числами (тому в реальному атомі гелію $\oa$ і $\ob$ додатково відрізняються ще й проєкцією спіна).

\section{Гамільтоніан та середня енергія}

Нерелятивістський гамільтоніан системи двох електронів у полі ядра складається з двох одноелектронних доданків $\h(1)$, $\h(2)$ (кінетична енергія кожного електрона плюс його притягання до ядра) та доданка міжелектронного відштовхування $1/r_{12}$:

\begin{equation*}
    \hat{H} = \ha + \hb + \V \quad E = \opbracket{\Phi(1, 2)}{\hat{H}}{\Phi(1, 2)}
\end{equation*}

Енергія системи в наближенні Хартрі--Фока --- це просто середнє значення оператора $\hat H$ на пробній хвильовій функції $\Phi(1,2)$. Підставимо явний вигляд $\Phi$ через $1/\sqrt2\,(\a-\b)$ (де $\a\equiv\oa\f\ob\s$, а $\b\equiv\ob\f\oa\s$ --- це два доданки визначника Слейтера) і розкриємо матричний елемент. Множник $1/\sqrt2$ з бра- та кет-вектора дає загальний коефіцієнт $1/2$, а сам добуток розкладається на чотири доданки виду $\langle \a|\dots|\a\rangle$, $\langle\a|\dots|\b\rangle$, $\langle\b|\dots|\a\rangle$, $\langle\b|\dots|\b\rangle$ для кожного з трьох операторів $\ha$, $\hb$, $\V$ окремо:

\hrulefill
\begin{multline*}
    E = \frac12 \opbracket{\cPhi}{\ha + \hb + \V}{\cPhi} = \\
    = \frac12\opbracket{\cPhi}{\ha}{\cPhi} + \\
    + \frac12\opbracket{\cPhi}{\hb}{\cPhi} + \\
    + \frac12\opbracket{\cPhi}{\V}{\cPhi} = \\[10pt]
    \hhr
    % for h1
    = \frac12\opbracket{\a}{\ha}{\a} + \frac12\opbracket{\b}{\ha}{\b} - \\
    - \frac12\opbracket{\a}{\ha}{\b} - \frac12\opbracket{\b}{\ha}{\a} + \\
    \hr
    % for h2
    + \frac12\opbracket{\a}{\hb}{\a} + \frac12\opbracket{\b}{\hb}{\b} - \\
    - \frac12\opbracket{\a}{\hb}{\b} - \frac12\opbracket{\b}{\hb}{\a} + \\
    \hr
    % for V
    + \frac12\opbracket{\a}{\V}{\a} + \frac12\opbracket{\b}{\V}{\b} - \\
    - \frac12\opbracket{\a}{\V}{\b} - \frac12\opbracket{\b}{\V}{\a} = \\
    \hhr
    % for h1
    = \frac12\opbracket{\oa\f}{\ha}{\oa\f} + \frac12\opbracket{\ob\f}{\ha}{\ob\f} + \\
    % for h2
    + \frac12\opbracket{\oa\s}{\hb}{\oa\s} + \frac12\opbracket{\ob\s}{\hb}{\ob\s} + \\
    \hr
    % for J
    + \frac12\opbracket{\a}{\V}{\a} + \frac12\opbracket{\b}{\V}{\b} - \\
    % for K
    - \frac12\opbracket{\a}{\V}{\b} - \frac12\opbracket{\b}{\V}{\a} =  \\
    \hhr
    % for h
    = \opbracket{\F}{\h}{\F} + \opbracket{\S}{\h}{\S}
    % for J
    + \frac12\opbracket{\F\S}{\V}{\F\S} + \frac12\opbracket{\S\F}{\V}{\S\F} -
    % for K
    \frac12\opbracket{\F\S}{\V}{\S\F} - \frac12\opbracket{\S\F}{\V}{\F\S}
\end{multline*}

\subsection*{Пояснення до кожного кроку розкриття}

\begin{itemize}[leftmargin=1.4em]
    \item \textbf{Крок 1 (доданки з $\ha$).} Оператор $\ha\equiv\h(1)$ діє лише на координату електрона 1 і \emph{не чіпає} координату електрона 2. Тому в матричному елементі $\opbracket{\a}{\ha}{\b}=\opbracket{\oa\f\ob\s}{\h(1)}{\ob\f\oa\s}$ інтегрування за координатами електрона 2 дає множник $\langle\ob\s|\oa\s\rangle=0$ в силу ортогональності орбіталей. Такі <<перехресні>> доданки

     $\opbracket{\a}{\ha}{\b}$ та $\opbracket{\b}{\ha}{\a}$ зникають --- це видно вже на наступному кроці, де вони просто відсутні. Доданки $\opbracket{\a}{\ha}{\a}$ та $\opbracket{\b}{\ha}{\b}$, навпаки, факторизуються: інтеграл за координатами електрона 2 дає $\langle\ob\s|\ob\s\rangle=1$ (нормування), і залишається чисто одноелектронний інтеграл за координатою електрона 1 --- звідси $\opbracket{\a}{\ha}{\a}=\opbracket{\oa\f}{\ha}{\oa\f}$ і аналогічно для другого доданка.
    \item \textbf{Крок 1 (доданки з $\hb$).} Абсолютно симетрично: $\hb\equiv\h(2)$ діє лише на координату електрона 2, перехресні доданки зникають через ортогональність $\langle\oa\f|\ob\f\rangle=0$, а діагональні --- факторизуються, даючи одноелектронні інтеграли за координатою електрона 2.
    \item \textbf{Доданки з $\V=1/r_{12}$.} Оператор кулонівського відштовхування діє одночасно на обидва електрони і \emph{не факторизується} --- він не є сумою одноелектронних операторів. Тому тут жоден із чотирьох доданків не зникає й не спрощується до одноелектронного вигляду: усі чотири --- це справжні двоелектронні (шестивимірні) інтеграли. Доданки $\opbracket{\a}{\V}{\a}$, $\opbracket{\b}{\V}{\b}$ (з однаковим порядком орбіталей у бра й кет) називають \emph{кулонівськими} внесками, а $\opbracket{\a}{\V}{\b}$, $\opbracket{\b}{\V}{\a}$ (з переставленими координатами) --- \emph{обмінними}: вони виникають виключно завдяки антисиметризації хвильової функції і не мають класичного аналога.
    \item \textbf{Об'єднання доданків з $\ha$ і $\hb$.} Після скорочення нулів у другому блоці залишаються по два одноелектронних доданки від $\ha$ і по два від $\hb$; попарно вони однакові (доданок від $\h(1)$ на орбіталі $\oa$ і доданок від $\h(2)$ на тій самій орбіталі $\oa$ дають однакове число, бо це лише перепозначення інтеграційної змінної). Це дозволяє в третьому блоці записати їх компактно через позначення $\F\equiv\oa$, $\S\equiv\ob$ (без прив'язки до конкретного електрона) і отримати $\opbracket{\F}{\h}{\F}+\opbracket{\S}{\h}{\S}$ --- суму одноелектронних (\emph{орбітальних}) енергій.
    \item \textbf{Останній блок.} Групуємо кулонівські доданки в $J$-внесок $\frac12\opbracket{\F\S}{\V}{\F\S}+\frac12\opbracket{\S\F}{\V}{\S\F}$ (за симетрією $1\leftrightarrow2$ ці два доданки рівні між собою) та обмінні --- в $K$-внесок з протилежним знаком.
\end{itemize}

\section{Кулонівський ($J$) та обмінний ($K$) інтеграли}

Отриманий вище результат зручно переписати, ввівши стандартні позначення
\[
    J \equiv \opbracket{\F\S}{\V}{\F\S} = \iint |\oa(1)|^2\,\frac{1}{r_{12}}\,|\ob(2)|^2\, d\vr_1\, d\vr_2,
 \]
\[
    K \equiv \opbracket{\F\S}{\V}{\S\F} = \iint \oa^*(1)\ob^*(2)\,\frac{1}{r_{12}}\,\ob(1)\oa(2)\, d\vr_1\, d\vr_2 .
\]

Фізичний зміст: $J$ --- це звичайна класична енергія електростатичного відштовхування двох <<хмаринок>> заряду з густинами $|\oa|^2$ і $|\ob|^2$. Обмінний інтеграл $K$ такого класичного тлумачення не має --- він виникає суто квантовомеханічно, з вимоги антисиметрії, і в чисельнику підінтегрального виразу орбіталі електронів 1 і 2 "переплетені" (кожна орбіталь залежить від координат обох електронів по черзі). У результаті повна енергія системи набуває вигляду

\begin{equation*}
    E = \ve_1 + \ve_2 + J - K = \sum_{i = 1}^{2}\ve_i + \frac12 \sum_{i = 1}^{2} \sum_{j = 1}^{2} (J_{ij} - K_{ij}), \quad J_{ii} = K_{ii} = 0 ,
\end{equation*}
де $\ve_i\equiv\opbracket{\vphi_i}{\h}{\vphi_i}$ --- орбітальна (одноелектронна) енергія, а подвійна сума в другій формі запису --- це той самий результат, записаний так, щоб узагальнення на довільну кількість електронів було очевидним: умова $J_{ii}=K_{ii}=0$ означає, що електрон не взаємодіє сам із собою (самодія відсутня, оскільки для $i=j$ кулонівський та обмінний доданки в цій формі тотожно скорочуються).

\section{Функціонал Фока і умова нормування}

Орбіталі $\vphi_i$ поки що довільні, і наша мета --- знайти саме ті з них, які мінімізують енергію $E$. Пряма мінімізація $E$ без обмежень безглузда: енергію можна зробити як завгодно малою, просто "розмазавши" орбіталь на нескінченність (втративши нормування). Тому мінімізацію проводять \emph{за умови збереження ортонормованості} орбіталей, $\langle\vphi_i|\vphi_j\rangle=\delta_{ij}$, застосовуючи метод невизначених множників Лагранжа $\ve_{ij}$:

\begin{equation*}
    F = E - \sum_{i = 1}^{2} \sum_{j = 1}^{2} (\ve_{ij} \bracket{\vphi_i}{\vphi_j} - \delta_{ij}) .
\end{equation*}

Множники $\ve_{ij}$ утворюють матрицю (в загальному випадку ермітову, $\ve_{ij}=\ve_{ji}^*$, що випливає з ермітовості умови нормування) --- саме вона згодом визначить структуру рівнянь ХФ. Далі варіюємо $F$ незалежно за $\vphi_1^*$ і $\vphi_2^*$ (орбіталі та їх комплексно спряжені трактуються як незалежні варіаційні параметри --- стандартний прийом), прирівнюючи варіацію до нуля.

\section{Варіювання функціонала}

Варіюючи кожен доданок енергії $E$ (одноелектронний, кулонівський і обмінний) та доданок з множниками Лагранжа за $\delta\vphi_1^*\equiv\delta\F$ і $\delta\vphi_2^*\equiv\delta\S$, отримуємо:

\begin{multline*}
    \delta F =  \opbracket{\delta\F}{\h}{\F} + \opbracket{\delta\S}{\h}{\S} + \\
    % for J
    + \frac12\opbracket{\delta\F\S}{\V}{\F\S} + \frac12\opbracket{\F\delta\S}{\V}{\F\S} +
    \frac12\opbracket{\delta\S\F}{\V}{\S\F} + \frac12\opbracket{\S\delta\F}{\V}{\S\F} - \\
    % for K
    - \frac12\opbracket{\delta\F\S}{\V}{\S\F} - \frac12\opbracket{\delta\S\F}{\V}{\F\S}
    - \frac12\opbracket{\F\delta\S}{\V}{\S\F} - \frac12\opbracket{\S\delta\F}{\V}{\F\S} - \\
    - \ve_{11} \bracket{\delta\F}{\F}
    - \ve_{12} \bracket{\delta\F}{\S}
    - \ve_{21} \bracket{\delta\S}{\F}
    - \ve_{22} \bracket{\delta\S}{\S} + \\
    + \text{c. c.} = 0 .
\end{multline*}

Кожен з двоелектронних ($J$- та $K$-) доданків енергії залежить одночасно від $\vphi_1$ \emph{і} $\vphi_2$, тому за правилом Лейбніца варіювання добутку дає по \emph{два} доданки --- один з варіацією по $\F$, інший по $\S$ (звідси подвоєння кожного з чотирьох доданків $J$ і $K$ до восьми). Оскільки нас цікавлять окремо коефіцієнти при незалежних варіаціях $\delta\F$ та $\delta\S$, наступний крок --- зібрати всі доданки, пропорційні $\delta\F$, в одне рівняння, а всі доданки, пропорційні $\delta\S$, --- в інше. Доданок "c.c." (комплексно спряжений вираз) автоматично дає нуль окремо, оскільки $\delta\F$ і $\delta\F^*$ незалежні, тому досить прирівняти нулю коефіцієнт при $\delta\F$ (і, окремо, при $\delta\S$):

\begin{align*}
    \left( \h + \frac12\opbracket{\dt\S}{\V}{\dt\S} + \frac12\opbracket{\S\dt}{\V}{\S\dt}
    - \frac12 \opbracket{\dt\S}{\V}{\S\dt} - \frac12\opbracket{\S\,.}{\V}{\dt\S}
    \right) \ket{\F} &= \ve_{11} \ket{\F} + \ve_{12}\ket{\S}
    \\[1.1\baselineskip]%
    \left( \h + \frac12\opbracket{.\,\F}{\V}{.\,\F} + \frac12\opbracket{\F\dt}{\V}{\F\dt}
    - \frac12 \opbracket{\dt\F}{\V}{\F\dt} - \frac12\opbracket{\F\dt}{\V}{\dt\F}
    \right) \ket{\S} &= \ve_{21} \ket{\F} + \ve_{22}\ket{\S}
\end{align*}

Тут крапка $\dt$ у брі/кеті позначає "порожнє місце", яке залишилось після того, як з нього "вийняли" $\delta\F$ (або $\delta\S$), перетворивши двоелектронний матричний елемент числа на одноелектронний \emph{оператор}, що діє на решту кет-вектора $\ket{\F}$ (відповідно $\ket{\S}$). Праворуч стоїть уся сума $\sum_j\ve_{ij}\ket{\vphi_j}$ --- саме тут у явному вигляді проявляється те, що рівняння для $\vphi_1$ і $\vphi_2$ пов'язані \emph{недіагональною} матрицею множників Лагранжа $\ve_{ij}$: доки ця матриця не діагональна, орбіталі $\F$ і $\S$ "перемішуються" одна з одною правою частиною.

\section{Канонічні рівняння Хартрі--Фока}

Матриця $\ve_{ij}$ ермітова, а тому її завжди можна діагоналізувати унітарним перетворенням орбіталей $\vphi_i\to\sum_k U_{ik}\vphi_k$. Таке перетворення не змінює самого визначника Слейтера $\Phi(1,2)$ (з точністю до несуттєвого загального фазового множника), бо визначник змінюється лише на множник $\det U$, а фізичні спостережувані --- енергія й густина --- від нього не залежать. Обираючи унітарне перетворення, що діагоналізує $\ve_{ij}$ (тобто робить $\ve_{12}=\ve_{21}=0$), переходимо до \emph{канонічного} базису орбіталей, в якому рівняння ХФ набувають найпростішого, "власного", вигляду --- права частина містить лише саму орбіталь, а не суміш з іншою:

\begin{align*}
    \left( \h + \opbracket{\dt\S}{\V}{\dt\S}
    - \opbracket{\dt\S}{\V}{\S\dt}
    \right) \ket{\F} &= \ve_{11} \ket{\F}
    \\[1.1\baselineskip]%
    \left( \h + \opbracket{\dt\F}{\V}{\dt\F}
    -  \opbracket{\dt\F}{\V}{\F\dt}
    \right) \ket{\S} &= \ve_{22} \ket{\S}
\end{align*}

Це і є \textbf{канонічні рівняння Хартрі--Фока} для системи двох електронів. Вираз у дужках --- оператор Фока $\hat f$ для відповідної орбіталі: сума одноелектронного гамільтоніана $\h$ та \emph{ефективного} потенціалу взаємодії з другим електроном, побудованого з його орбіталі. Кожне рівняння --- звичайне рівняння на власні значення $\hat f\,\ket{\vphi_i}=\ve_i\ket{\vphi_i}$, але сам оператор $\hat f$ залежить від орбіталей \emph{обох} електронів --- звідси необхідність розв'язувати систему самоузгоджено (self-consistent field): задавшись пробними орбіталями, будують оператор Фока, розв'язують рівняння, отримують нові орбіталі, і так до збіжності.

\section{Кулонівський та обмінний оператори}

Останнім кроком явно виділимо в операторі Фока два доданки як окремі оператори, що діють на орбіталь $\ket{\F}$ (аналогічно --- на $\ket{\S}$).

\subsection*{Кулонівський оператор}

Кулонівський оператор $\hat{J}_{\S}$ --- це \emph{локальний} (мультиплікативний) оператор: у кожній точці простору $\vr_1$ він множить хвильову функцію на класичний електростатичний потенціал, створений густиною заряду $|\ob(2)|^2$ другого електрона, усередненою (проінтегрованою) за всіма положеннями електрона 2:

\begin{equation*}
   \hat{J}_{\S}\ket{\F} = \opbracket{\dt\S}{\V}{\dt\S} \ket{\F} = \int \frac{\dt\ob\s\dt\ob\s}{r_{12}} d\s \cdot \oa\f =  \int \frac{\dt\ob\s\oa\f\ob\s}{r_{12}} d\s
\end{equation*}

Крапки в $\bra{\dt\S}\,\V\,\ket{\dt\S}$ вказують, що координата електрона 1 залишається "вільною" (не проінтегрованою) --- саме вона й перетворює цей вираз з числа на оператор, який далі діє на $\ket{\F}$. Інтеграл береться лише за координатою $\vr_2$ електрона 2; результат --- функція від $\vr_1$, яка потім домножується на $\oa(1)$.

\subsection*{Кулонівський інтеграл}

Якщо додатково "закрити" залишену вільною координату 1, домноживши зліва на $\bra{\F}$ і проінтегрувавши, отримаємо вже не оператор, а \emph{число} --- середнє значення кулонівської енергії відштовхування електронних хмаринок $\oa$ і $\ob$:

\begin{equation*}
   \bra{\F}\hat{J}_{\S}\ket{\F} = \bra{\F}\opbracket{\dt\S}{\V}{\dt\S} \ket{\F} = \iint \frac{\oa\f\ob\s\oa\f\ob\s}{r_{12}} d\s d\f
\end{equation*}

Це шестивимірний інтеграл (по $\vr_1$ і $\vr_2$), симетричний відносно перестановки електронів; підінтегральна густина $|\oa(1)|^2\,|\ob(2)|^2$ --- добуток двох "класичних" густин ймовірності, тому $J$ завжди дійсне й додатне число.

\begin{figure}[h!]
\centering
\begin{tikzpicture}[scale=0.85]
	% "розмазаний" електрон 2 -- густина |phi_2(r_2)|^2
	\shade[inner color=blue!65, outer color=white] (0,0) circle (2.2);
	\draw[blue!40] (0,0) circle (2.2);
	\fill (0,0) circle (1.1pt);
	\node[blue!70!black] at (0,-2.6) {\footnotesize $|\varphi_{\mathbf{2}}(\vec r_2)|^2$ --- <<розмазаний>> електрон 2};
	% точковий (локальний) електрон 1
	\fill[ball color=red!70] (55:3.2) circle (0.13) coordinate (E1);
	\node[red!70!black, right] at (55:3.55) {\footnotesize $\varphi_{\mathbf{1}}(\vec r_1)$ --- локальний електрон 1};
	% взаємодія
	\draw[dashed, thick, gray!70!black] (0,0) -- (E1) node[midway, above, sloped] {\footnotesize $1/r_{12}$};
\end{tikzpicture}
\caption{Кулонівський оператор $\hat J_2$: електрон 2 <<розмазаний>> у неперервну густину ймовірності $|\varphi_2(\vec r_2)|^2$ (проінтегрований, тому й не має конкретної точки), а електрон 1 залишається локальним -- у кожній його точці $\vec r_1$ оператор просто домножує $\varphi_1(\vec r_1)$ на число: класичний електростатичний потенціал, створений усією хмарою заряду $|\varphi_2|^2$. Саме тому $\hat J_2$ -- локальний (мультиплікативний) оператор.}
\label{fig:coulomb-operator}
\end{figure}

\subsection*{Обмінний оператор}

Обмінний оператор $\hat{K}_{\S}$, на відміну від кулонівського, \emph{нелокальний}: значення $(\hat{K}_{\S}\vphi_1)(\vr_1)$ у точці $\vr_1$ залежить від значення самої орбіталі $\vphi_1$ у \emph{всіх} інших точках простору (під інтегралом стоїть $\oa(2)$, а не $\oa(1)$) --- це прямий наслідок антисиметризації хвильової функції, суто квантовий ефект без класичного аналога:

\begin{equation*}
    \hat{K}_{\S}\ket{\F} = \opbracket{\dt\S}{\V}{\S\dt}\ket{\F} = \int \frac{\dt\ob\s \ob\f\dt}{r_{12}} d\s \cdot \oa\f = \int \frac{\dt\ob\s \ob\f\oa\s}{r_{12}} d\s
\end{equation*}

Порівняно з $\hat{J}_{\S}$ тут орбіталі $\oa$ і $\ob$ під інтегралом "обмінялися" аргументами: замість $|\ob(2)|^2\oa(1)$ маємо $\ob(1)\ob^*(2)\oa(2)$ --- звідси й назва "обмінний".

\subsection*{Обмінний інтеграл}

Аналогічно кулонівському, "закриваючи" вільну координату бра-вектором $\bra{\F}$, отримуємо обмінний інтеграл:

\begin{equation*}
    \bra{\F}\hat{K}_{\S}\ket{\F} = \bra{\F}\opbracket{\dt\S}{\V}{\S\dt}\ket{\F} = \iint \frac{\oa\f\ob\s \ob\f\oa\s}{r_{12}} d\s d\f
\end{equation*}

На відміну від кулонівського інтеграла, тут під знаком інтеграла не густини ймовірності, а добутки $\oa^*(1)\ob(1)$ та $\ob^*(2)\oa(2)$ --- так звані \emph{густини переходу} між орбіталями $\oa$ і $\ob$. Незважаючи на це, можна показати, що для дійсних орбіталей $K\ge0$, а знак "$-K$" в енергії ($E=\ve_1+\ve_2+J-K$) означає, що обмінна взаємодія завжди \emph{понижує} повну енергію системи щодо чисто кулонівського результату --- це відображення того, що антисиметрична (фермі-статистична) хвильова функція автоматично "розводить" електрони з паралельними спінами в просторі, зменшуючи їхнє середнє відштовхування (обмінна діра Фермі).

\begin{figure}[h!]
\centering
\begin{tikzpicture}[scale=0.85]
	% "розмазана" перехідна густина phi_2^*(2) phi_1(2) -- вже не густина ймовірності, тому колір -- сумішевий
	\shade[inner color=violet!55, outer color=white] (0,0) circle (2.2);
	\draw[violet!35] (0,0) circle (2.2);
	\fill (0,0) circle (1.1pt);
	\node[violet!60!black] at (0,-2.6) {\footnotesize $\varphi_{\mathbf{2}}^*(\vec r_2)\,\varphi_{\mathbf{1}}(\vec r_2)$ --- <<перехідна густина>>};
	% точковий електрон 1: після дії К його "ярлик" підмінюється на 2
	\fill[ball color=red!70] (55:3.2) circle (0.13) coordinate (E1);
	\node[right] at (55:3.55) {\footnotesize $\varphi_{\mathbf{1}}(\vec r_1)\ \xrightarrow{\ \hat K_2\ }\ \textcolor{blue!70!black}{\varphi_{\mathbf{2}}(\vec r_1)}$};
	% взаємодія
	\draw[dashed, thick, gray!70!black] (0,0) -- (E1) node[midway, above, sloped] {\footnotesize $1/r_{12}$};
\end{tikzpicture}
\caption{Обмінний оператор $\hat K_2$: <<хмара>> електрона 2 тепер побудована не з густини $|\varphi_2|^2$, а з \emph{перехідної} густини $\varphi_2^*(\vec r_2)\varphi_1(\vec r_2)$, яка одночасно <<пам'ятає>> про обидві орбіталі -- суто квантовий, некласичний об'єкт. Оператор $\hat K_2$ нелокальний: подіявши на $\varphi_1(\vec r_1)$ у точці $\vec r_1$, він повертає результат, пропорційний $\varphi_2(\vec r_1)$ -- орбіталь електрона 1 <<підмінюється>> орбіталлю електрона 2, що й відображає їх фізичну нерозрізненість.}
\label{fig:exchange-operator}
\end{figure}

\end{document}
